\section{Conclusion}

\label{sec:conclusion}
% ============================================================================

Base is not a backend framework. It is the runtime for sovereign applications.
The distinction matters. A framework assumes an environment---a cloud provider,
a database server, a container orchestrator. A runtime \emph{is} the
environment. Base carries its own database, its own auth, its own function
execution, its own sync protocol. It runs anywhere a Go binary runs.

The design principles are:

\begin{enumerate}[leftmargin=1.4em]
    \item \textbf{Local first}: The device is the source of truth. The network
    is an optimization, not a requirement.
    \item \textbf{Encrypted by default}: Data at rest is encrypted. Field-level
    encryption is available for sensitive data. Keys are user-controlled.
    \item \textbf{Sync, don't migrate}: CRDT-based replication means data flows
    between devices without schema migrations or conflict resolution dialogs.
    \item \textbf{Chain-optional}: Lux integration provides key policy,
    identity, and threshold controls for applications that need them. Those
    that do not need them pay no cost.
    \item \textbf{One binary}: No external dependencies. No infrastructure
    prerequisites. Download, run, build.
\end{enumerate}

The personal computer gave individuals sovereignty over computation. The
smartphone gave individuals sovereignty over communication. Base gives
individuals sovereignty over their data. That is the thesis. The 10,400
encrypted writes per second are just the proof.

% ============================================================================
% REFERENCES
% ============================================================================
\begin{thebibliography}{99}

\bibitem{hanzobase}
Z.~Kelling.
\newblock Hanzo Base: An open-source backend runtime with embedded database,
authentication, and {JavaScript} plugin system.
\newblock \emph{Hanzo AI Research}, December 2022.

\bibitem{luxwhitepaper}
{Lux Partners}.
\newblock Lux: A scalable multi-consensus blockchain with post-quantum
cryptography.
\newblock \emph{Lux Network Technical Report}, 2024.

\bibitem{luxtrust}
Z.~Kelling.
\newblock The Lux Trust Kernel: Threshold cryptography for decentralized key
management.
\newblock \emph{Lux Network Technical Report}, 2025.

\bibitem{luxapp}
Z.~Kelling.
\newblock Lux Application Architecture: Sovereign applications on
multi-consensus infrastructure.
\newblock \emph{Lux Network Technical Report}, 2025.

\bibitem{hanzoplatform}
M.~Chen, D.~Wei, and Z.~Kelling.
\newblock Hanzo Platform: {PaaS} for {AI}-native application deployment.
\newblock \emph{Hanzo AI Research}, December 2021.

\bibitem{kleppmann2019}
M.~Kleppmann, A.~Wiggins, P.~van Hardenberg, and M.~McGranaghan.
\newblock Local-first software: You own your data, in spite of the cloud.
\newblock In \emph{Onward!}, ACM, 2019.

\bibitem{automerge}
M.~Kleppmann and A.~R.~Beresford.
\newblock A conflict-free replicated {JSON} datatype.
\newblock \emph{IEEE TPDS}, 28(10):2733--2746, 2017.

\bibitem{yjs}
K.~Jahns.
\newblock Yjs: A {CRDT} framework for shared editing.
\newblock \url{https://github.com/yjs/yjs}, 2019.

\bibitem{kulkarni2014hlc}
S.~Kulkarni, M.~Demirbas, D.~Madeppa, B.~Avva, and M.~Leone.
\newblock Logical physical clocks and consistent snapshots in globally
distributed databases.
\newblock In \emph{OPODIS}, 2014.

\bibitem{shapiro2011crdt}
M.~Shapiro, N.~Pregui\c{c}a, C.~Baquero, and M.~Zawirski.
\newblock Conflict-free replicated data types.
\newblock In \emph{SSS}, Springer, 2011.

\bibitem{cryptdb}
R.~A.~Popa, C.~M.~S.~Redfield, N.~Zeldovich, and H.~Balakrishnan.
\newblock {CryptDB}: Protecting confidentiality with encrypted query processing.
\newblock In \emph{SOSP}, ACM, 2011.

\bibitem{mylar}
R.~A.~Popa, E.~Stark, J.~Helfer, S.~Valdez, N.~Zeldovich,
M.~F.~Kaashoek, and H.~Balakrishnan.
\newblock Building web applications on top of encrypted data using {Mylar}.
\newblock In \emph{NSDI}, USENIX, 2014.

\bibitem{pocketbase2022}
G.~Kostadinov.
\newblock {PocketBase}: Open-source backend in 1 file.
\newblock \url{https://pocketbase.io}, 2022.

\end{thebibliography}

